\begin{ExerciseBlock}
\ExerciseHeading{07}{Reibung mit seitlicher Strömung}{forces_lateral_fluid_drag}
Eine Objekt bewegt sich auf einer geraden Fahrspur (z.B. Schiene) mit einer Geschwindigkeit $v_F$. Das Gas oder die Flüssigkeit, die das Objekt umgeben, strömt mit einer Geschwindigkeit $v_W$ senkrecht zur Fahrtrichtung. \\
Zwischen dem Objekt und dem umgebenden Medium kommt es zur Reibung. Die Reibungskraft $\vec{F}_\mathrm R$ hat eine Komponente senkrecht zur Fahrtrichtung $F_\mathrm R^\perp$, die keine Auswirkung auf die Bewegung des Objektes hat, und eine Komponente parallel zur Fahrtrichtung $F_\mathrm R^{\parallel}$, die das Objekt abbremst.

\begin{center}
	\CachedTikzFigure{content/resources/chapters/forces/tikz/lateral_fluid_drag.tex}
\end{center}

\begin{enumerate}
	\item Berechnen Sie einen Ausdruck die Kraftvektor $\vec{F}_\mathrm R$ für den Fall einer laminaren Reibung (Stokes-Reibung).
	\[ \vec{F}_\mathrm R = -b \vec{v} \]
	Geben Sie die Komponente $F_\mathrm R^{\parallel}$ an.

	\item Berechnen Sie einen Ausdruck die Kraftvektor $\vec{F}_\mathrm R$ für den Fall einer Newton'schen Reibung.
	\[ \vec{F}_\mathrm R = - k \; v^2 \vec{e}_v \]
	Geben Sie die Komponente $F_\mathrm R^{\parallel}$ an.

	\item Welchen Unterschied beobachten Sie zwischen Aufgabenteil a) und b). Welchen Einfluss hat die Geschwindigkeit der seitlichen Strömung $v_W$ jeweils auf $F_\mathrm R^{\parallel}$?

	\item Nehmen Sie an, das Objekt (z.B. ein Fahrradfahrer) bewege sich mit $v_F=15\,$km/h und sei von Luft umgeben. Um wieviel Prozent erhöht sich die Reibungskraft $F_\mathrm R^{\parallel}$ (Newton'sche Reibung) bei einem Seitenwind von $v_W=10 \,$km/h gegenüber einer Situation mit Windstille?

\end{enumerate}
\end{ExerciseBlock}
\UseExerciseSolution{forces_lateral_fluid_drag}
